\begin{center}
{\zihao{3}\bfseries Abstract}
\end{center}

Battery packs are representative complex electromechanical products in new-energy vehicle manufacturing. Their assembly process involves structural integration, electrical connection, thermal management, and safety testing, and therefore imposes stringent requirements on process coordination, inspection accuracy, and manufacturing consistency. Although current production lines have introduced robots, machine vision, and digital traceability, assembly and inspection are still often organized in a fragmented manner. Process knowledge, equipment capability, and quality data are not yet sufficiently integrated, which limits flexibility and closed-loop decision-making in high-mix, low-volume manufacturing scenarios.

This paper takes the battery pack as a concrete product and examines its assembly tasks, process chain, inspection nodes, digital means, key intelligent technologies, and a new assembly mode. First, the product structure of the battery pack is used to analyze major assembly tasks and critical inspection points. Second, the roles and coordination mechanisms of robots, human--robot collaboration, machine vision, AI, knowledge graphs, and large language model agents are discussed in the assembly environment. Finally, a knowledge-driven human--robot collaborative mode for battery pack intelligent assembly and online inspection is proposed. In this mode, robots, human--robot collaboration, and machine vision constitute the execution layer, while process knowledge graphs, rule reasoning, and large language model agents constitute the cognitive support layer, with emphasis on explicit representation of process constraints, closed-loop feedback from online inspection, and assisted decision-making for abnormal conditions.

The analysis shows that the key to intelligent battery pack assembly does not lie in introducing a single advanced algorithm, but in integrating process knowledge, equipment capability, inspection results, and human decisions into a unified system. The proposed mode has potential to improve changeover efficiency, abnormal response, and quality traceability. However, its practical deployment is still constrained by knowledge maintenance cost, industrial data governance, model trustworthiness, and safety certification. Overall, the paper develops a coherent discussion framework covering product-oriented process analysis, system design, key technologies, application prospects, and limitations.

\vspace{1em}
\noindent\textbf{Keywords:} battery pack; intelligent assembly; online inspection; knowledge graph; large language model agent
